% Part: first-order-logic
% Chapter: completeness
% Section: construction-of-model

\documentclass[../../../include/open-logic-section]{subfiles}

\begin{document}

\iftag{FOL}
      {\olfileid{fol}{com}{mod}}
      {\olfileid{pl}{com}{mod}}

\olsection{إنشاء نموذج}

\begin{explain}
\iftag{FOL}{نؤخر النظر في $\eq$، ونقصد الآن إلى إرضاء كل مجموعة
  متسقة~$\Gamma$ من ال!!{sentence}s خالية من~$\eq$. نمد~$\Gamma$
  أولًا إلى مجموعة متسقة و!!{complete} ومشبعة~$\Gamma^*$.
  ثم نعرّف النموذج~$\Struct{M(\Gamma^*)}$ تعريفًا يسيرًا: مجاله
  الحدود المغلقة في~$\Lang{L'}$، وكل !!{constant} يُفسَّر بنفسه.
  ويعم هذا الحكم الحدود المغلقة، فنضمن $\Value{t}{M(\Gamma^*)} = t$
  لكل حد مغلق~$t$. وأما ال!!{predicate}s فنعين امتداداتها بحيث
  تصدق ال!!{sentence} الذرية في $\Struct{M(\Gamma^*)}$ إذا وفقط
  إذا كانت في~$\Gamma^*$. وبذلك تصدق ال!!{sentence}s الذرية كلها
  في~$\Gamma^*$ داخل $\Struct{M(\Gamma^*)}$. وأما غير الذرية،
  فيكفل صدقها اتساق $\Gamma^*$ واكتمالها وإشباعها.}{حان تعريف
  !!a{valuation} ترضي كل $!A \in \Gamma$. نأخذ، بلمّة ليندنبوم،
  مجموعة متسقة كاملة $\Gamma^* \supseteq \Gamma$، ثم نعيّن التقييم
  بما يقع من ال!!{propositional variable}s في~$\Gamma^*$؛ وليكن
  هذا التقييم~$\pAssign v(\Gamma^*)$.}
\end{explain}

\iftag{FOL}{%
\begin{defn}[نموذج الحدود]
\ollabel{defn:termmodel}
إذا كانت $\Gamma^*$ مجموعة !!a{complete} ومتسقة ومشبعة من ال!!{sentence}s
في لغة~$\Lang L$، سمينا \emph{نموذج الحدود}~$\Struct M(\Gamma^*)$ الخاص
بـ$\Gamma^*$ ال!!{structure} التي نحددها بهذه البنود:
\begin{enumerate}
\item نجعل ال!!{domain}~$\Domain{M(\Gamma^*)}$ مجموعة الحدود المغلقة
  كلها في~$\Lang L$.
\item نأخذ !!a{constant} $c$؛ فتفسيره $c$ بعينه:
  $\Assign{c}{M(\Gamma^*)} = c$.
\item نفسر ال!!{function}~$f$ بالدالة التي ترسل الحدود المغلقة
  $t_1$، \dots، $t_n$ إلى الحد المغلق
  $f(t_1, \dots, t_n)$؛ أي:
\[
\Assign{f}{M(\Gamma^*)}(t_1, \dots, t_n) = f(t_1,\dots, t_n)
\]
\item إذا كان $R$ !!{predicate} ذا $n$ مواضع، فإن
  \[
  \tuple{t_1, \dots,
  t_n} \in \Assign{R}{M(\Gamma^*)} \text{ إذا وفقط إذا } \Atom{R}{t_1, \dots,
    t_n} \in \Gamma^*.
  \]
\end{enumerate}
\end{defn}
}{%
\begin{defn}
\ollabel{defn:termmodel}
متى كانت $\Gamma^*$ مجموعة متسقة كاملة من ال!!{formula}s، عرّفنا التقييم الآتي:
\[
\pAssign v(\Gamma^*)(p) = \begin{cases}
  \True & \text{إذا كان $p \in \Gamma^*$}\\
  \False & \text{إذا كان $p \notin \Gamma^*$}
  \end{cases}
\]
\end{defn}
}

\iftag{FOL}{%
نثبت أولًا أن التفسير يفي بما وعدنا به، أي $\Value{t}{M(\Gamma^*)} = t$.

\begin{lem}
\ollabel{lem:val-in-termmodel} ليكن $\Struct M(\Gamma^*)$ نموذج الحدود
في \olref{defn:termmodel}؛ عندئذ $\Value{t}{M(\Gamma^*)} = t$ لكل حد
مغلق~$t$.
\end{lem}
\begin{proof}
 نستقرئ تركيب $t$. ففي الأساس يكون $t$ !!a{constant}، وقيمته معلومة
 من تعريف نموذج الحدود. ولخطوة الاستقراء خذ حدودًا مغلقة
 $t_1, \ldots, t_n$ تحقق $\Value{t_i}{M(\Gamma^*)} = t_i$، وليكن
 $f$ !!{function} له $n$ حجج. عندئذ نحسب:
\begin{align*}
\Value{f(t_1,\ldots,t_n)}{M(\Gamma^*)} &= \Assign{f}{M(\Gamma^*)}(\Value{t_1}
 {M(\Gamma^*)},
\ldots, \Value{t_n}{M(\Gamma^*)}) \\
&= \Assign{f}{M(\Gamma^*)}(t_1, \dots, t_n) \\
&= f(t_1,\dots, t_n),
\end{align*}
ومن ثم، بالاستقراء، يصح ذلك لكل حد مغلق~$t$.
\end{proof}
}

\iftag{FOL}{%
\begin{explain}
  \iftag{prvEx}{يجوز أن ترضي !!^a{structure}~$\Struct{M}$
    !!{sentence} وجودية~$\lexists[x][!A(x)]$، ولا ترضي أي
    حالة~$!A(t)$ منها.}{}
  \iftag{prvAll}{ويجوز أن ترضي !!^a{structure}~$\Struct{M}$ كل
    حالة~$!A(t)$ من !!{sentence} كلية~$\lforall[x][!A(x)]$،
    ولا ترضي~$\lforall[x][!A(x)]$ نفسها.}{} ووجه ذلك أن
  !!{element} من~$\Domain{M}$ قد لا يكون قيمة لأي حد مغلق؛ فليس
  لازمًا أن تكون $\Struct{M}$ مغطاة. ولهذا احتجنا إلى إسنادات
  المتغيرات في تعريف الإرضاء. فأما نموذج الحدود~$\Struct{M(\Gamma^*)}$
  فمغطى، فلا حاجة فيه إلى ذلك الطريق؛ وهذا ما تفصله القضية الآتية.
\end{explain}

\begin{prop}
\ollabel{prop:quant-termmodel}
ليكن $\Struct M(\Gamma^*)$ نموذج الحدود في \olref{defn:termmodel}.
\begin{tagenumerate}{prvEx,prvAll}
\tagitem{prvEx}{$\Sat{M(\Gamma^*)}{\lexists[x][!A(x)]}$ إذا وفقط إذا
  كان $\Sat{M(\Gamma^*)}{!A(t)}$ لحد مغلق واحد~$t$ على الأقل.}{}
\tagitem{prvAll}{$\Sat{M(\Gamma^*)}{\lforall[x][!A(x)]}$ إذا وفقط إذا
  كان $\Sat{M(\Gamma^*)}{!A(t)}$ لكل حد مغلق~$t$.}{}
\end{tagenumerate}
\end{prop}

\begin{proof}
\begin{tagenumerate}{prvEx,prvAll}
\tagitem{prvEx}{%
  \iftag{probEx}{تمرين.}{نستعمل \olref[syn][ass]{prop:sat-quant}: فصدق
    $\Sat{M(\Gamma^*)}{\lexists[x][!A(x)]}$ يكافئ وجود إسناد
    متغيرات~$s$ يحقق $\Sat{M(\Gamma^*)}{!A(x)}[s]$.
    ومجال $\Domain{M(\Gamma^*)}$ هو الحدود المغلقة في~$\Lang{L}$؛
    فذلك يكافئ وجود حد مغلق~$t$ مع إسناد يحقق
    $s(x) = t$ و$\Sat{M(\Gamma^*)}{!A(x)}[s]$. ثم تفيد
    \olref[fol][syn][ext]{prop:ext-formulas},
    يتحقق $\Sat{M(\Gamma^*)}{!A(x)}[s]$ إذا وفقط إذا تحقق
    $\Sat{M(\Gamma^*)}{!A(t)}[s]$، حيث $s(x) = t$. وبحسب
    \olref[fol][syn][ass]{prop:sentence-sat-true}، يتحقق
    $\Sat{M(\Gamma^*)}{!A(t)}[s]$ إذا وفقط إذا تحقق
    $\Sat{M(\Gamma^*)}{!A(t)}$، لأن $!A(t)$ جملة.}}{}
\tagitem{prvAll}{%
  \iftag{probAll}{تمرين.}{تقضي \olref[syn][ass]{prop:sat-quant} بأن
    $\Sat{M(\Gamma^*)}{\lforall[x][!A(x)]}$ يكافئ، لكل إسناد متغيرات
    $s$، تحقق $\Sat{M(\Gamma^*)}{!A(x)}[s]$. وعناصر
    $\Domain{M(\Gamma^*)}$ هي الحدود المغلقة في~$\Lang{L}$؛ فيمكن
    اختيار إسناد يحقق $s(x) = t$ لأي حد مغلق~$t$، وبالعكس تكون
    $s(x)$ في كل إسناد حدًا مغلقًا~$t$. وتفيد
    \olref[fol][syn][ext]{prop:ext-formulas}، يتحقق
    $\Sat{M(\Gamma^*)}{!A(x)}[s]$ إذا وفقط إذا تحقق
    $\Sat{M(\Gamma^*)}{!A(t)}[s]$، حيث $s(x) = t$. وبحسب
    \olref[fol][syn][ass]{prop:sentence-sat-true}، يتحقق
    $\Sat{M(\Gamma^*)}{!A(t)}[s]$ إذا وفقط إذا تحقق
    $\Sat{M(\Gamma^*)}{!A(t)}$، لأن $!A(t)$ جملة.}}{}
\end{tagenumerate}
\end{proof}
}{}

\tagprob[FOL]{probEx,probAll}
\begin{prob}
أكمل برهان \olref[fol][com][mod]{prop:quant-termmodel}.
\end{prob}
\tagendprob

\begin{lem}[لمّة الصدق]
\ollabel{lem:truth} \iftag{FOL}{افترض أن $!A$ لا تحتوي~$\eq$. عندئذ}{}
$\iftag{FOL}{\Sat{M(\Gamma^*)}}{\pSat{v(\Gamma^*)}}{!A}$ إذا وفقط إذا
كان $!A \in \Gamma^*$.
\end{lem}

\begin{proof}
نجري الاستقراء على $!A$، ونثبت وجهي التكافؤ في كل حالة معًا.
\begin{enumerate}
\tagitem{prvFalse}{\indcase{!A}{\lfalse}
  {$\iftag{FOL}{\Sat/{M(\Gamma^*)}{\lfalse}}{\pSat/{v(\Gamma^*)}{\lfalse}}$
    هو حكم تعريف الإرضاء؛ ويقابله $\lfalse \notin \Gamma^*$، إذ
    اتساق $\Gamma^*$ يمنع اشتمالها على الباطل.}}{}

\tagitem{prvTrue}{\indcase{!A}{\ltrue}
  {$\iftag{FOL}{\Sat{M(\Gamma^*)}}{\pSat{v(\Gamma^*)}}{\ltrue}$
    يقتضيه تعريف الإرضاء؛ وكذلك $\ltrue \in \Gamma^*$، لاجتماع
    الاتساق وكون $\Gamma^*$ !!{complete} مع $\Gamma^* \Proves \ltrue$.}}{}

\tagitem{FOL}{\indcase{!A}{R(t_1, \dots, t_n)}
  {$\Sat{M(\Gamma^*)}{\Atom{R}{t_1, \dots, t_n}}$ إذا وفقط إذا كان $\tuple{t_1,
      \dots, t_n} \in \Assign{R}{M(\Gamma^*)}$ (بحسب تعريف الإرضاء)
    إذا وفقط إذا كان $R(t_1, \dots, t_n) \in \Gamma^*$ (بحسب إنشاء
    $\Struct M(\Gamma^*)$).}}{\indcase{!A}{p}{$\pSat{v(\Gamma^*)}{p}$
    إذا وفقط إذا كان $\pAssign v(\Gamma^*)(p) = \True$ (بحسب تعريف
    الإرضاء)، وذلك إذا وفقط إذا كان $p \in \Gamma^*$ (بحسب إنشاء
    $\pAssign v(\Gamma^*)$).}}

\tagitem{prvNot}{%
  \indcase{!A}{\lnot !B}
    {$\iftag{FOL}{\Sat{M(\Gamma^*)}}{\pSat{v(\Gamma^*)}}{\indfrm}$ إذا وفقط إذا تحقق
    $\iftag{FOL}{\Sat/{M(\Gamma^*)}{!B}}{\pSat/{v(\Gamma^*)}{!B}}$
    بتعريف الإرضاء. ثم ينقلنا فرض الاستقراء إلى أن
    $\iftag{FOL}{\Sat/{M(\Gamma^*)}{!B}}{\pSat/{v(\Gamma^*)}{!B}}$ إذا
    وفقط إذا كان $!B \notin \Gamma^*$. والاتساق مع كون $\Gamma^*$
    !!{complete} يقتضيان أن $!B \notin \Gamma^*$ يكافئ
    $\lnot !B \in \Gamma^*$.}}{}

\tagitem{prvAnd}{%
\iftag{probAnd}{%
  \indcase!{!A}{!B \land !C}{}}{%
  \indcase{!A}{!B \land
    !C}{$\iftag{FOL}{\Sat{M(\Gamma^*)}}{\pSat{v(\Gamma^*)}}{\indfrm}$
    إذا وفقط إذا كان كل من
    $\iftag{FOL}{\Sat{M(\Gamma^*)}}{\pSat{v(\Gamma^*)}}{!B}$ و
    $\iftag{FOL}{\Sat{M(\Gamma^*)}}{\pSat{v(\Gamma^*)}}{!C}$ متحققًا
    بتعريف الإرضاء؛ ويحوّل فرض الاستقراء هذا إلى اجتماع
    $!B \in \Gamma^*$ و$!C \in \Gamma^*$. ثم تفيد
    \olref[ccs]{prop:ccs}\olref[ccs]{prop:ccs-and}، يتحقق ذلك إذا وفقط
    إذا كان $(!B \land !C) \in \Gamma^*$.}}}{}

\tagitem{prvOr}{%
\iftag{probOr}{%
  \indcase!{!A}{!B \lor !C}{}}{%
  \indcase{!A}{!B \lor !C}
    {$\iftag{FOL}{\Sat{M(\Gamma^*)}}{\pSat{v(\Gamma^*)}}{\indfrm}$
    إذا وفقط إذا تحقق
    $\iftag{FOL}{\Sat{M(\Gamma^*)}}{\pSat{v(\Gamma^*)}}{!B}$ أو
    $\iftag{FOL}{\Sat{M(\Gamma^*)}}{\pSat{v(\Gamma^*)}}{!C}$ بتعريف
    الإرضاء؛ وبفرض الاستقراء يكافئ ذلك $!B \in \Gamma^*$ أو
    $!C \in \Gamma^*$. وهذان البديلان يكافئان
    $(!B \lor !C) \in \Gamma^*$، بمقتضى
    \olref[ccs]{prop:ccs}\olref[ccs]{prop:ccs-or}.}}}{}

\tagitem{prvIf}{%
\iftag{probIf}{%
  \indcase!{!A}{!B \lif !C}{}}{%
  \indcase{!A}{!B \lif !C}
    {$\iftag{FOL}{\Sat{M(\Gamma^*)}}{\pSat{v(\Gamma^*)}}{\indfrm}$
    إذا وفقط إذا تحقق
    $\iftag{FOL}{\Sat/{M(\Gamma^*)}}{\pSat/{v(\Gamma^*)}}{!B}$ أو
    $\iftag{FOL}{\Sat{M (\Gamma^*)}}{\pSat{v(\Gamma^*)}}{!C}$ بتعريف
    الإرضاء؛ وبتطبيق فرض الاستقراء نجد $!B \notin \Gamma^*$ أو
    $!C \in \Gamma^*$. وذلك مكافئ لـ
    $(!B \lif !C) \in \Gamma^*$، بمقتضى
    \olref[ccs]{prop:ccs}\olref[ccs]{prop:ccs-if}.}}}{}

\iftag{FOL}{%
\tagitem{prvAll}{%
  \iftag{probAll}{%
    \indcase!{!A}{\lforall[x][!B(x)]}{}}{%
    \indcase{!A}{\lforall[x][!B(x)]}{$\Sat{M(\Gamma^*)}{\indfrm}$ إذا
      وفقط إذا كان $\Sat{M(\Gamma^*)}{!B(t)}$ لكل حد مغلق~$t$
      (\olref{prop:quant-termmodel}). وهذا، بفرض الاستقراء، يكافئ
      $!B(t) \in \Gamma^*$ لكل حد مغلق~$t$؛ ثم ترده
      \olref[hen]{prop:saturated-instances} إلى
      $\lforall[x][!B(x)] \in \Gamma^*$.}}}{}

\tagitem{prvEx}{%
  \iftag{probEx}{%
    \indcase!{!A}{\lexists[x][!B(x)]}{}}{%
    \indcase{!A}{\lexists[x][!B(x)]}{$\Sat{M(\Gamma^*)}{\indfrm}$ إذا
      وفقط إذا كان $\Sat{M(\Gamma^*)}{!B(t)}$ لحد مغلق واحد~$t$ على
      الأقل (\olref{prop:quant-termmodel}). وينقلنا فرض الاستقراء إلى
      الشرط المكافئ $!B(t) \in \Gamma^*$ لحد مغلق~$t$ واحد على
      الأقل. وتفيد \olref[hen]{prop:saturated-instances} أن هذا
      يكافئ $\lexists[x][!B(x)] \in \Gamma^*$.}}}{}
}{}
\end{enumerate}
\end{proof}


\tagprob[FOL]{probOr,probAnd,probIf,probEx,probAll}

\begin{prob}
أكمل برهان \olref[fol][com][mod]{lem:truth}.
\end{prob}

\tagendprob

\tagprob[notFOL]{probOr,probAnd,probIf}

\begin{prob}
أكمل برهان \olref[pl][com][mod]{lem:truth}.
\end{prob}

\tagendprob

\end{document}
